\section{贪污贿赂罪 - 贪污罪、挪用公款罪}
\begin{enumerate}
  \item \zd{贪污}。A + B。本质是财产犯罪，特殊的侵占、盗窃、诈骗罪，\zx{B = B1 + B2 + B3。缺一个 fallback}。
        \begin{enumerate}
          \item \zd{把构成要件的他人改成公家，就很好理解}。如他人/公家所有、他人/公家占有 \yjt 自己占有、所有。
          \item B1 \yjt 行为主体。\zx{国家工作人员}。
          \item B2 \yjt 行为对象。\bs{现实拥有、已然存在的公家财物}。

          \item B3 \yjt 行为方式。从事\zd{公务}过程中利用职务便利 \yjt 职务有一定\zx{行政职责、行政管理性}。
                \begin{enumerate}
                  \item \al{甲掌握钥匙，乙掌握密码。乙\xhx{捡到}甲的钥匙打开钱箱。利用的是自己有密码的便利。乙盗窃型}。
                  \item \al{乙骗到甲的钥匙打开钱箱。乙盗窃型}。
                  \item \al{甲偷看乙的密码打开钱箱。利用了自己有钥匙的便利。甲盗窃型}。
                  \item \al{甲乙共谋}。
                \end{enumerate}
          \item A \yjt 实行行为。
                \begin{enumerate}
                  \item 侵占/侵吞。
                        \begin{enumerate}
                          \item 转移所有不要求职务便利，占有已经体现了职务便利。\al{出纳保管钥匙，为规避犯罪，不用钥匙开锁，而是锤子砸开}。
                          \item \zx{既遂。行使所有权}。
                        \end{enumerate}
                  \item 盗窃/窃取。
                        \begin{enumerate}
                          \item 转移占有要求职务便利。\al{甲保管钱箱钥匙和密码，取走现金 \yjt 不要求职务便利，甲侵占型}。
                          \item \zx{既遂。建立自己的占有}。
                        \end{enumerate}
                  \item 诈骗/骗取。一般是\zx{有职权（明显发挥了作用）}的领导，才能利用职务便利诈骗。道理同上。
                        \begin{enumerate}
                          \item \al{甲有权\xhx{管理工资}事宜，伪造工资表领工资，会计乙只实施了形式审查。甲诈骗型}。
                          \item \al{\xhx{科员}甲为多报销，\xhx{伪造报销单}，结果财务有一次就被骗了。fallback，甲普通诈骗}。
                        \end{enumerate}
                \end{enumerate}
          \item 数额。等号原理。\zx{被害人损失数额 = 犯罪人所得数额}。
                \al{国有公司甲领导，给乙老板出借 100W，还钱时候甲让乙老板打到个人账户 60W，帮做平账}。
          \item \bs{容易遗漏 A，只看到了 B1、B2、B3 \dui}。
                \al{\xhx{甲（普通小偷）}要盗窃国有企业仓库，找到保安经理乙（国家工作人员），请乙行个方便。甲盗窃，乙盗窃帮助犯}。
          \item 共犯与\bs{身份}。
                \begin{enumerate}
                  \item 非甲帮助，\zx{国乙实行}。\al{民政局乙想贪污，让甲提供假发票。乙贪污，甲贪污帮助犯}。
                  \item \zx{非甲实行}，国乙帮助。
                        \al{甲为骗领补助，\xhx{民政局}乙\xhx{帮其成为低保户}身份，甲向\xhx{社保局}申领 1W 补助。甲诈骗，乙诈骗帮助犯}。
                  \item \zx{非甲实行，国乙实行}。
                        \begin{enumerate}
                          \item \al{投保人甲和国有企业经理乙内外勾结，相互勾结，骗取保险公司保险金}。
                          \item \al{村民甲虚构危房补助材料，民政局长乙\xhx{明知，仍批给}甲 \yjt 处分行为即为实行行为。乙贪污}。
                        \end{enumerate}
                \end{enumerate}
          \item 奇葩案例。\al{国甲曾因公务为公司垫付各种费用 5W，
                  因票据超期无法报销。
                  \xhx{为挽回损失}，和私乙虚构劳务合同并虚开发票，找公司财物套取劳务费 5W。啥都不构成}。
        \end{enumerate}
  \item 挪用公款。\zx{挪出行为 + 归“个人”使用}。
        \begin{enumerate}
          \item \zx{实行行为。只有挪，用叫客观处罚条件}。
                \al{官员甲挪用资金贩毒 \yjt 如果是挪 + 用，挪用评价一次，用于贩毒评价一次，\xhx{重复评价}}。
          \item 公款。\bs{只能是钱}。\al{国甲以公司\xhx{在售的商品房}偿还自己的债务}。
          \item 归个人。\bs{欠你人情 \hz 你得好处}。\al{工商局长甲\xhx{以单位名义}将公款借给某公司使用、装修办公大楼 \cuo}。
          \item 3 个使用方式。\zx{只要求客观具备，不需要主客观一致}。\al{本来\xhx{想去赌}后面\xhx{改成炒股}}。
                \begin{enumerate}
                  \item \zx{非法活动 \yjt 3W}。不要求达 3 个月。
                  \item \zx{营利活动 \yjt 5W}。\bs{钱生钱}。\al{做生意、炒股}。
                  \item \zx{一般活动 \yjt 5W + 超 3 个月未还}。\sjt \zx{构成其他罪的，数罪并罚}。\al{挪用公款还收人贿赂}。
                \end{enumerate}
          \item 数额认定。\zx{\yi 包容评价。\er 实际发生额}。
                \begin{enumerate}
                  \item \al{挪用 7W，2.5W \xhx{用于非法}，4.5W \xhx{用于营利} \yjt 营利\xhx{包容评价}非法，认定 7W 进行营利活动}。
                  \item \al{乙\xhx{怂恿}国甲挪用 300W 给自己\xhx{投资}，称 2 月后归还并承诺获利平分。
                          后乙\xhx{用 100W 投资}，200W 购房。
                          \xhx{两个月后}乙将 300W 还于甲，甲归还给公司。挪用数额，甲乙均是 100W}。
                \end{enumerate}
          \item vs 贪污。要求非法占有目的（不想还），挪用想还。\zx{可转化为贪污}（产生不想还的目的）。
                \al{甲挪用单位办公经费 70W 为自己买房，两周后，采取\xhx{销毁账本}目的，使 70W 中的 50W 难以在账务上反映出来}。
          \item \zx{逃避处罚，携款潜逃原则不定贪污}。
          \item 共犯与\bs{身份}。
                \begin{enumerate}
                  \item \al{\xhx{国甲挪用}公款和非乙贩毒。甲挪用公款、贩卖毒品，乙贩卖毒品}。
                  \item \al{非乙为贩毒\xhx{指使}国甲挪用公款。甲挪用公款、贩卖毒品，乙挪用公款教唆犯、贩卖毒品。数罪并罚}。
                  \item \al{乙\xhx{谎称}为了看病\xhx{唆使}甲挪用公款，称 2 周后归还，乙实际是去贩毒，但 2 周后确实还了。甲挪用公款客观阶段实行犯，但最终不构成。乙挪用公款教唆犯、贩卖毒品，数罪并罚}。
                \end{enumerate}
        \end{enumerate}
\end{enumerate}
